\documentclass[11pt]{article}
\usepackage{lemme}
\usepackage{tabularx,array,colortbl}
\renewcommand\tabularxcolumn[1]{m{#1}}
\renewcommand\lemmemeta{Formation commando X-ENS · Topologie · 1}
\renewcommand\lemmeurl{lemme-ryq6.vercel.app/\#commando}
\renewcommand{\leq}{\leqslant}\renewcommand{\geq}{\geqslant}
\newcommand{\R}{\mathbb{R}}
\newcommand{\K}{\mathbb{K}}
\begin{document}

\lemmetitle{Formation commando X-ENS · Résumé de cours}{Topologie : les généralités}

\lemmenotions{Dans tout le résumé, $(E,d)$ est un espace métrique et $\K = \R$ ou $\mathbb{C}$.}

\lemmesection{1. Normes et distances}

\begin{ldef}[Norme]
Une norme sur un $\K$-espace vectoriel $E$ est une application $x\mapsto\|x\|$ à valeurs dans $\R_+$ telle que
\[\|x\| = 0\iff x = 0,\qquad \|\lambda x\| = |\lambda|\,\|x\|,\qquad \|x+y\|\leq\|x\|+\|y\|.\]
\end{ldef}

Dans $\R^n$, les trois normes à connaître sont $\|x\|_1 = \sum|x_i|$, $\|x\|_2 = \sqrt{\sum x_i^2}$ et $\|x\|_\infty = \max|x_i|$. Pour les fonctions bornées sur un ensemble $X$ : $\|f\|_\infty = \sup_X|f|$, la norme de la convergence uniforme.

\begin{ldef}[Distance]
Une distance sur un ensemble $E$ vérifie : $d(x,y) = 0\iff x = y$, \ $d(x,y) = d(y,x)$, \ $d(x,z)\leq d(x,y)+d(y,z)$.
\end{ldef}

\begin{itemize}
\item Une norme donne une distance : $d(x,y) = \|x-y\|$. C'est celle qu'on utilise par défaut.
\item La \textbf{distance discrète} vaut $0$ si $x = y$ et $1$ sinon. Elle sert de contre-exemple.
\item Diamètre : $\delta(A) = \sup_{x,y\in A}d(x,y)$. Une partie est \textbf{bornée} si son diamètre est fini, c'est-à-dire si elle tient dans une boule.
\item Distance d'un point à une partie non vide : $d(x,A) = \inf_{y\in A}d(x,y)$. Attention, $d(A,B) = 0$ n'implique pas $A = B$.
\end{itemize}

\lemmesection{2. Boules, ouverts, fermés}

$B(a,r) = \{x : d(a,x)<r\}$ est la boule ouverte, $B_f(a,r) = \{x : d(a,x)\leq r\}$ la boule fermée, $S(a,r) = \{x : d(a,x) = r\}$ la sphère.

\begin{ldef}[Ouvert, fermé]
$\Omega$ est \textbf{ouvert} si autour de chacun de ses points, il contient une boule : $\forall x\in\Omega,\ \exists\rho>0,\ B(x,\rho)\subset\Omega$. \\
$F$ est \textbf{fermé} si son complémentaire est ouvert.
\end{ldef}

\begin{lthm}[Stabilité]
\begin{itemize}
\item $\varnothing$ et $E$ sont ouverts et fermés.
\item Une réunion \emph{quelconque} d'ouverts est ouverte ; une intersection \emph{finie} d'ouverts est ouverte.
\item Une intersection \emph{quelconque} de fermés est fermée ; une réunion \emph{finie} de fermés est fermée.
\end{itemize}
\end{lthm}

\begin{lnote}[Attention]
Une intersection infinie d'ouverts peut ne pas être ouverte : $\bigcap_{n\geq1}\left]-\frac1n,\frac1n\right[ = \{0\}$.
\end{lnote}

Une boule ouverte est un ouvert, une boule fermée est un fermé. $V$ est un \textbf{voisinage} de $x$ s'il contient un ouvert qui contient $x$.

\lemmesection{3. Adhérence, intérieur, frontière}

\begin{ldef}[Adhérence et intérieur]
L'\textbf{adhérence} $\overline A$ est le plus petit fermé qui contient $A$. \\
L'\textbf{intérieur} $\mathring A$ est le plus grand ouvert contenu dans $A$. \\
La \textbf{frontière} est $\operatorname{Fr}(A) = \overline A\setminus\mathring A$.
\end{ldef}

\begin{lthm}[Caractérisations]
\[x\in\overline A\iff\forall\varepsilon>0,\ B(x,\varepsilon)\cap A\neq\varnothing\iff d(x,A) = 0.\]
\[x\in\mathring A\iff\exists\varepsilon>0,\ B(x,\varepsilon)\subset A.\]
$A$ est fermé si et seulement si $A = \overline A$ ; ouvert si et seulement si $A = \mathring A$.
\end{lthm}

\begin{itemize}
\item $A$ est \textbf{dense} dans $E$ si $\overline A = E$. Par exemple, $\mathbb{Q}$ et $\R\setminus\mathbb{Q}$ sont denses dans $\R$.
\item Dans un espace normé, $\overline{B(0,1)} = B_f(0,1)$. Dans un espace métrique quelconque, on a seulement l'inclusion $\subset$.
\item $a$ est un \textbf{point isolé} de $A$ si une boule autour de $a$ ne contient aucun autre point de $A$. $a$ est un \textbf{point d'accumulation} de $A$ si toute boule autour de $a$ contient un point de $A$ autre que $a$.
\end{itemize}

\begin{ldef}[Topologie induite]
Sur une partie $A$ munie de la distance restreinte, les ouverts sont les $\Omega\cap A$ et les fermés les $F\cap A$, où $\Omega$ est ouvert et $F$ fermé dans $E$.
\end{ldef}

\lemmesection{4. Continuité}

\begin{ldef}[Continuité en un point]
$f : (E,d)\to(E',d')$ est continue en $a$ si
\[\forall\varepsilon>0,\ \exists\alpha>0,\ \forall x\in E,\quad d(a,x)<\alpha\implies d'(f(a),f(x))<\varepsilon.\]
\end{ldef}

\begin{lthm}[Continuité globale]
$f$ est continue sur $E$ si et seulement si l'image réciproque de tout ouvert est un ouvert, si et seulement si l'image réciproque de tout fermé est un fermé.
\end{lthm}

\begin{itemize}
\item La composée de fonctions continues est continue. Dans un espace normé, la somme et le produit par un scalaire sont continus.
\item L'image directe d'un ouvert par une fonction continue n'est pas forcément ouverte.
\item Un \textbf{homéomorphisme} est une bijection continue dont la réciproque est continue.
\end{itemize}

\begin{ldef}[Continuité uniforme et fonctions lipschitziennes]
$f$ est \textbf{uniformément continue} si le même $\alpha$ convient pour tous les points :
\[\forall\varepsilon>0,\ \exists\alpha>0,\ \forall x,y\in E,\quad d(x,y)<\alpha\implies d'(f(x),f(y))<\varepsilon.\]
$f$ est \textbf{$k$-lipschitzienne} si $d'(f(x),f(y))\leq k\,d(x,y)$ pour tous $x,y$.
\end{ldef}

\begin{lthm}[La chaîne à retenir]
Lipschitzienne $\implies$ uniformément continue $\implies$ continue. \\
Les réciproques sont fausses : $x\mapsto\sqrt x$ est uniformément continue sur $\R_+$ sans être lipschitzienne, et $x\mapsto\frac1x$ est continue sur $]0,1]$ sans être uniformément continue.
\end{lthm}

\lemmesection{5. Normes et distances équivalentes}

Deux normes $N_1$ et $N_2$ sont \textbf{équivalentes} s'il existe $a,b>0$ avec $aN_1\leq N_2\leq bN_1$. Elles définissent alors les mêmes ouverts, les mêmes fermés, les mêmes fonctions continues. En dimension finie, toutes les normes sont équivalentes.

\lemmesection{6. Produits et limites}

\begin{itemize}
\item Sur $E = E_1\times\dots\times E_n$, la distance produit est $d(x,y) = \max_i d_i(x_i,y_i)$. Ses boules sont des produits de boules.
\item $f = (f_1,\dots,f_n)$ est continue si et seulement si chaque composante $f_i$ l'est.
\item Une fonction continue en chaque variable séparément n'est pas forcément continue : $\frac{xy}{x^2+y^2}$, prolongée par $0$ en $(0,0)$.
\item La distance $d : E\times E\to\R$ est continue (et même $2$-lipschitzienne).
\item Les limites se composent, passent aux sommes et aux produits, et respectent les inégalités larges, comme dans $\R$.
\end{itemize}

\newpage
\lemmesection{Les formules à connaître}

\renewcommand{\arraystretch}{1.55}\arrayrulecolor{line}
\begin{tabularx}{\linewidth}{>{\sffamily\small\color{muted}\raggedright\arraybackslash}m{4.4cm} X}
Inégalité triangulaire & $d(x,z)\leq d(x,y)+d(y,z)$ \\ \hline
Inégalité renversée & $|d(x,z)-d(y,z)|\leq d(x,y)$ \\ \hline
Normes usuelles sur $\R^n$ & $\|x\|_1 = \sum|x_i|$, \quad $\|x\|_2 = \sqrt{\sum x_i^2}$, \quad $\|x\|_\infty = \max|x_i|$ \\ \hline
Comparaison sur $\R^n$ & $\|x\|_\infty\leq\|x\|_2\leq\|x\|_1\leq n\,\|x\|_\infty$ \\ \hline
Diamètre & $\delta(A) = \sup_{x,y\in A}d(x,y)$ \\ \hline
Distance à une partie & $d(x,A) = \inf_{y\in A}d(x,y)$ \\ \hline
Lipschitz de $d(\cdot,A)$ & $|d(x,A)-d(y,A)|\leq d(x,y)$ \\ \hline
Boules & $B(a,r) = \{d(a,x)<r\}$, \quad $B_f(a,r) = \{d(a,x)\leq r\}$ \\ \hline
Ouvert & $\forall x\in\Omega,\ \exists\rho>0,\ B(x,\rho)\subset\Omega$ \\ \hline
Adhérence & $x\in\overline A\iff\forall\varepsilon>0,\ B(x,\varepsilon)\cap A\neq\varnothing\iff d(x,A) = 0$ \\ \hline
Intérieur & $x\in\mathring A\iff\exists\varepsilon>0,\ B(x,\varepsilon)\subset A$ \\ \hline
Complémentaires & $\mathring A = E\setminus\overline{E\setminus A}$, \quad $\overline A = E\setminus(E\setminus A)^{\circ}$ \\ \hline
Frontière & $\operatorname{Fr}(A) = \overline A\setminus\mathring A$ \\ \hline
Continuité & $\forall\varepsilon>0,\ \exists\alpha>0,\ d(a,x)<\alpha\implies d'(f(a),f(x))<\varepsilon$ \\ \hline
Continuité uniforme & $\forall\varepsilon>0,\ \exists\alpha>0,\ \forall x,y,\ d(x,y)<\alpha\implies d'(f(x),f(y))<\varepsilon$ \\ \hline
Lipschitzienne & $d'(f(x),f(y))\leq k\,d(x,y)$ \\ \hline
Normes équivalentes & $a\,N_1\leq N_2\leq b\,N_1$ \\ \hline
Distance produit & $d(x,y) = \max_i d_i(x_i,y_i)$ \\ \hline
Max de deux réels & $\max(a,b) = \frac12\big(a+b+|a-b|\big)$ \\ \hline
\end{tabularx}

\end{document}
